\documentclass[11pt]{article}
\usepackage{amsmath,amssymb,amsthm,hyperref}
\newtheorem{theorem}{Theorem}
\newtheorem{corollary}{Corollary}
\newcommand{\E}{\mathbb E}
\title{Repeated comparison columns force a larger support}
\author{Asymptotic research note; independently reviewed}
\date{30 September 2026}
\begin{document}
\maketitle
Let $p_1,\ldots,p_m:[K]\to[0,1]$ be nondecreasing, and give the
ordered rows arbitrary positive weights. Suppose
\[
 \E\prod_{j\in S}p_j=\frac1{|S|+1}
 \qquad(S\subseteq[m]).
\]
This is the necessary comparison-tensor condition for a distinguished
die with $K$ faces in a permutation-fair family of $m+1$ dice.

\begin{theorem}\label{main}
Suppose at least $r\ge2$ columns equal a common function $a$.
There is an integer $R\in[r,m]$ such that $a$ has the uniform-interval
moments through degree $R$, and either
\begin{enumerate}
\item $R=m$, so $a$ itself is a scalar quadrature of degree $m$; or
\item $R<m$, and
\[
 K\ge
 \left\lfloor\frac R2\right\rfloor(m-R)
 +\left\lceil\frac R2\right\rceil+1.
\]
\end{enumerate}
If $a$ is strictly increasing, one can take $R$ to be exactly the
number of columns equal to $a$. The theorem also permits arbitrary
plateaus in every column.
\end{theorem}
\begin{proof}
For every column outside the original equal group put $e_j=p_j-a$.
The first two mixed moments, using two copies of $a$, give
\[
 \E e_j=\E ae_j=0,\qquad \E e_i e_j=0\quad(i\ne j).
\]
Call $e_j$ flat if $\E(e_j\mid a)=0$. On every level block $B$ of
$a$, all columns remain nondecreasing, and a flat error has mean zero.
For a flat $e_i$ and any other error $e_j$,
\[
 \E(e_i e_j\mid B)=\operatorname{Cov}(e_i,e_j\mid B)\ge0.
\]
Their unconditional inner product is zero. Positivity of the row
weights therefore makes this conditional covariance zero on each block.
The equality case of the monotone covariance inequality says that
at least one function is constant on the block: its covariance is
\[
 \frac12\sum_{s,t\in B}w_sw_t
 (e_i(s)-e_i(t))(e_j(s)-e_j(t)),
\]
and its first--last term is positive if both functions vary.

Consequently, wherever a flat error varies, every other error is
constant. Moreover at most one flat error varies on a block; a flat
error that is constant on a block is zero there. Replacing all flat
columns by $a$ preserves every mixed moment: on each block a product
either remains unchanged or has one varying mean-zero error times
constant other factors.

Let $R$ be the size of the enlarged equal group, and let
$k=m-R$ errors remain. If $k=0$, the tensor is scalar after replacement,
and all moments through degree $m$ follow. Otherwise every remaining
error has nonzero conditional mean given $a$.
Choosing distinct columns from the enlarged group gives
\[
 \E a^s=\frac1{s+1}\quad(0\le s\le R),\qquad
 \E a^s e_j=0\quad(0\le s\le R-1),
\]
and, for distinct remaining errors,
\[
 \E a^s e_i e_j=0\quad(0\le s\le R-2).
\]
The signed measure obtained by grouping $e_j$ over the distinct values
of $a$ is nonzero and annihilates polynomials of degree at most $R-1$.
It must therefore be supported at at least $R+1$ distinct values,
by the Vandermonde determinant. In particular $a$ has at least $R+1$
values, and every remaining $e_j$ is nonzero on at least $R+1$ of
its level sets.

Put $d=\lfloor(R-2)/2\rfloor$. The $k$ subspaces
\[
 U_j=\operatorname{span}\{e_j,ae_j,\ldots,a^d e_j\}
\]
are mutually orthogonal, each of dimension $d+1=\lfloor R/2\rfloor$.
Their dimensions follow by evaluating any vanishing polynomial
combination on the more than $d$ distinct levels supporting $e_j$.
They are also orthogonal to
\[
 P=\operatorname{span}\{1,a,\ldots,a^{R-1-d}\},
 \qquad \dim P=R-d=\lceil R/2\rceil+1.
\]
Indeed all inner products involved have total power at most $R-2$
for two error spaces and at most $R-1$ for $P$ and an error space.
These orthogonal spaces lie in the $K$-dimensional row space, giving
the claimed bound.

If $a$ is strictly increasing, conditioning on $a$ fixes the row,
so a flat error is identically zero. Only columns already equal to
$a$ are added, as asserted.
\end{proof}

\begin{corollary}[Equal-weight asymptotic consequences]
For equal row weights there is an absolute constant $c>0$ such that
every tensor as above with $r\ge2$ identical columns satisfies
\[
 K\ge c\,m r.
\]
If the entries are rational, the same $R$ from Theorem~\ref{main}
also satisfies
\[
 K\ge 2^{2\lfloor R/2\rfloor}.
\]
In particular, a rational comparison tensor with $K\le m^C$ for a fixed
constant $C$ can
contain only $O(\log m)$ identical columns. When a repeated class
has size of order $\log m$, the new dimension estimate adds an
$\Omega(m\log m)$ restriction in the regime $R\le m/2$.
\end{corollary}
\begin{proof}
Equal-weight interval quadrature of degree $R$ requires $K\ge c_0R^2$
for an absolute $c_0>0$, by the classical Chebyshev quadrature lower
bound. If $R\ge m/2$, this implies $K\ge(c_0/2)mR$.
If $R<m/2$, the dimension bound gives
$K\ge\lfloor R/2\rfloor(m-R)\ge mR/6$.
Since $R\ge r$, the first conclusion follows.
The arithmetic conclusion uses the independently proved optimized
2-adic bound in
\path{../../independent_directions/optimized_two_adic_filter.tex}.
\end{proof}

The classical real quadrature input is recorded, for example, in
Kuijlaars, \emph{The minimal number of nodes in Chebyshev type
quadrature formulas} (1993),
\href{https://doi.org/10.1016/0019-3577(93)90007-L}{doi:10.1016/0019-3577(93)90007-L}.
This note does not prove a universal quadratic per-die lower bound:
models with many distinct comparison columns remain outside its
effective range. Its conclusion applies to actual dice whenever
their comparison columns satisfy the stated repetition hypothesis.
\end{document}
